\chapter{INTRODUCTION}
Hyperspectral images (HSI) capture electromagnetic reflectance across hundreds of narrow spectral bands per pixel, encoding the physical and chemical properties of Earth's surface materials at a level of detail that conventional RGB imagery simply cannot achieve. This spectral richness makes HSI indispensable in agriculture, environmental monitoring, and defense applications. Yet the same abundance of information that gives HSI its analytical power also introduces a fundamental computational burden: the well-documented curse of dimensionality inflates training costs for machine learning and deep learning models, while the scarcity of expert-labeled pixel samples compounds the problem by severely limiting model generalization.

This project investigates the performance gap between supervised and semi-supervised deep learning architectures---CNN, ViT, and Self-Training---under these twin constraints of high dimensionality and labeled data scarcity. Two contributions distinguish it from prior work. First, Curriculum Learning is integrated into the training pipeline to stabilize optimization across all model variants. Second, a multithreaded Graphical User Interface (GUI) is developed so that users can load external data and obtain predictions in real time, without touching a single line of code.

The report is structured as follows. Chapter 2 surveys the relevant literature and situates the project within the current state of the art. Chapter 3 details the system design and methodology, covering data preprocessing through model architectures. Chapter 4 presents comparative experimental results across multiple performance metrics, and Chapter 5 discusses findings and outlines directions for future work.

\section{Purpose}
The primary objective of this project is to quantitatively characterize the performance differences between supervised and semi-supervised learning in the classification of high-dimensional hyperspectral data. This overarching goal breaks down into the following sub-objectives:
\begin{itemize}
    \item To benchmark model performance on three standard hyperspectral datasets widely adopted in the literature: Indian Pines, Pavia University, and Salinas.
    \item To assess how preprocessing steps---Gaussian filtering for noise suppression and Principal Component Analysis (PCA) for dimensionality reduction---affect classification accuracy.
    \item To incorporate Curriculum Learning into each model to improve training stability and convergence behavior.
    \item To compare models objectively using Accuracy, Precision, Recall, F1-Score, and Kappa metrics.
    \item To deliver a Python-based multithreaded GUI that exposes the trained models to end users, enabling real-time prediction without requiring direct code interaction.
\end{itemize}

\section{Preliminary Investigation}
Surveying the hyperspectral classification literature reveals a clear trend: traditional machine learning methods are steadily being displaced by deep learning approaches. A recurring limitation across existing systems, however, is their dependence on substantial quantities of labeled data---a resource that is both costly and slow to produce, since pixel annotation requires domain expertise and substantial manual effort.

This project targets that gap directly. Semi-supervised algorithms are designed to exploit a small labeled set alongside a large pool of unlabeled data, and the experiments here test whether that design advantage translates into competitive classification performance under realistic labeling budgets. A second gap is more practical in nature: most research models exist only as scripts, inaccessible to non-technical stakeholders. The multithreaded GUI introduced in this project bridges that divide, turning experimental code into an application that accepts user-supplied data and returns instantaneous predictions.
